% chapters/08/section-05.tex
\section{Program Perilaku Mengejar}
\label{sec:program-chase}

Bagian ini menyatukan seluruh unsur menjadi program yang memuat beberapa NPC
yang mengejar seorang pemain pada ruang terbuka.

\subsection{Kode Lengkap}

\kodepython{python/ch08_chase.py}{Perilaku mengejar pada beberapa NPC sekaligus}{lst:ch08-chase}

\subsection{Membedah Keputusan NPC}

Seluruh keputusan NPC terletak pada satu fungsi, yaitu \code{Npc.update}.
Susunannya mengikuti tiga langkah pada Bagian~\ref{sec:bergerak-sendiri}.

\begin{description}
  \item[Persepsi] Dua baris pertama menghitung vektor menuju pemain dan
        jaraknya. Keduanya dihitung dari satu pengurangan yang
        sama, sehingga tidak ada perhitungan yang berulang.
  \item[Keputusan] Rangkaian pemeriksaan jarak menentukan kecepatan. Urutan
        pemeriksaan penting: jarak henti diperiksa lebih dahulu, baru jari-jari
        persepsi, sebab jarak henti selalu lebih kecil daripada jari-jari.
  \item[Tindakan] Posisi diperbarui dengan kecepatan tersebut, dibatasi agar
        tidak keluar jendela, lalu \code{rect} diselaraskan sebagaimana biasa.
\end{description}

Kelas \code{Npc} tidak menerima arah dari luar, berbeda dari
kelas \code{Player} pada bab ini dan pada bab-bab sebelumnya. Itulah perbedaan
yang paling mendasar antara benda yang diperintah dan benda yang memutuskan.

\subsection{Mengamati Keputusan melalui Gambar}

Program ini menggambar empat hal untuk memperlihatkan keputusan NPC, sesuai
uraian pada Bagian~\ref{sec:persepsi}: lingkaran jari-jari persepsi, lingkaran
jarak henti, garis arah menuju pemain, dan warna NPC yang berubah menurut
keputusannya.

Ketika pemain bergerak menjauh, warna NPC berubah menjadi abu-abu dan lingkaran
persepsinya tidak lagi memuat pemain. Ketika pemain mendekat, NPC berubah warna
dan mulai bergerak. Garis arahnya selalu menunjuk ke pemain, dan panjangnya
tidak mempengaruhi kecepatan NPC---itulah wujud indikator \ind{b} pada layar.

\subsection{Percobaan yang Disarankan}

\begin{table}[htbp]
\centering
\caption{Percobaan pada program perilaku mengejar}
\label{tab:percobaan-ch08}
{\setlength{\tabcolsep}{4pt}\footnotesize%
\begin{tabularx}{\textwidth}{@{}L{5.6cm} Y@{}}
\toprule
\headrow \textbf{Yang diubah dan caranya} & \textbf{Yang sebaiknya Anda amati} \\
\midrule
Jari-jari persepsi, melalui pilihan \code{--radius}. &
Pada jari-jari kecil, pemain dapat berlalu-lalang dekat NPC tanpa dikejar. Pada
jari-jari besar, NPC mulai bergerak bahkan ketika pemain berada jauh. \\
\addlinespace
Kecepatan NPC, melalui pilihan \code{--speed}. &
Pada kecepatan yang jauh lebih rendah daripada pemain, NPC tidak pernah
menangkap pemain. Pada kecepatan yang lebih tinggi, pemain selalu tertangkap,
dan NPC berhenti pada jarak hentinya. \\
\addlinespace
Jarak henti, dengan mengubah \code{DEFAULT_STOP} menjadi nilai yang sangat
kecil, misalnya 2. &
NPC bergetar di sekitar pemain, karena perpindahan setiap bingkainya lebih
besar daripada jarak henti. \\
\addlinespace
Cara menghitung kecepatan, dengan menjalankan program memakai pilihan
\code{--naive}. &
Kecepatan NPC berubah menurut jarak: lambat ketika jauh, cepat ketika dekat.
Bandingkan dengan keluaran pengukuran pada terminal. \\
\addlinespace
Banyaknya NPC, melalui pilihan \code{--npc}. &
Dengan banyak NPC, gerombolan akan mengelilingi pemain. Perhatikan pula beban
perhitungan ketika jumlahnya bertambah besar. \\
\bottomrule
\end{tabularx}}
\end{table}

\begin{catatabox}
Pilihan \code{--naive} sengaja disediakan untuk memperlihatkan cara yang salah.
Jalankan program dengan dan tanpa pilihan itu, lalu bandingkan kecepatan
terukur pada terminal. Perbedaan itulah bukti indikator \ind{b} pada bab ini.
\end{catatabox}
